\item Un estudiante mide los siguientes datos en un experimento de calorimetría diseñado para determinar el calor específico del aluminio:
\begin{itemize}
    \item Temperatura inicial de agua y calorímetro: $70^\circ\text{C}$
    \item Masa de agua: $0.400\text{ kg}$
    \item Masa del calorímetro: $0.040\text{ kg}$
    \item Calor específico del calorímetro: $0.63\text{ kJ/kg}\cdot^\circ\text{C}$
    \item Temperatura inicial del aluminio: $27.0^\circ\text{C}$
    \item Masa de aluminio: $0.200\text{ kg}$
    \item Temperatura final de mezcla: $66.3^\circ\text{C}$
\end{itemize}
(a) Use estos datos para determinar el calor específico del aluminio. (b) Explique si su resultado está dentro del 15\% del valor dado en la tabla 7.1.
